% Deutsche Parallelfassung zu state/admissibility.tex

\chapter{Zulässigkeitszustände}
\label{chap:admissibility-states}

Zulässigkeit ist keine zusätzliche geometrische Zustandsschicht. Sie ist die
zweckqualifizierte Bewertung am Ende der in den vorangehenden Kapiteln
entwickelten Abhängigkeitskette. Sie fragt, ob abgeleitete Größen innerhalb
einer anwendbaren betrieblichen Umgrenzung bleiben; sie ändert weder das
laufende Alignment noch macht sie aus einem Rechenergebnis eine Entscheidung.

\section{Zulässigkeitszustand}

\begin{definition}[Zulässigkeitszustand]
\label{def:admissibility-state}
Ein Zulässigkeitszustand ist ein Laufzeitbewertungszustand, der beschreibt, ob
betriebliche Größen innerhalb erlaubter Umgrenzungen bleiben.
\end{definition}

Zulässigkeitszustände können Grenzen für effektive Beschleunigung, Ruck,
Überhöhungsfehlbetrag, Rollentwicklung, Geschwindigkeit und
Realisierungsqualität umfassen. Diese Größen werden aus
Laufzeit-Eisenbahnzustand und dynamischer Bewertung abgeleitet; sie werden
nicht als Alignment-Identität dauerhaft gespeichert.

\section{Betriebliche Umgrenzungen}

\begin{definition}[Zustand der betrieblichen Umgrenzung]
\label{def:operational-envelope-state}
Ein Zustand der betrieblichen Umgrenzung beschreibt zulässige Gebiete im
Laufzeitzustandsraum.
\end{definition}

Betriebliche Umgrenzungen werden üblicherweise durch Normen,
Infrastrukturbeschränkungen und Betriebspolitik definiert. In AIM wirken sie
als Nebenbedingungen auf abgeleitete Laufzeitgrößen und nicht als Ersatz für
intrinsische Geometrie oder Eisenbahnzustand.

\section{Schichtposition}

Zulässigkeit wird nach pose2-/pose3-Rekonstruktion, metrischer Realisierung,
Frame-Konstruktion und Ableitung geschwindigkeitsqualifizierter Größen
bewertet. Regel, Version, Zweck, anwendbares Objekt und anwendbarer Zustand,
Gültigkeitszeit und Autorität müssen angegeben sein. Die Bewertung bleibt dem
Eisenbahnzustand nachgelagert und definiert ihn nicht neu.

Zulässigkeit auf Fahrzeugebene wird relativ zum Eisenbahnzustand bewertet und
darf nicht mit der Eisenbahnzustandsschicht selbst verwechselt werden.

\section{Verbindung zu AIM}

\begin{summary}
Zulässigkeitszustände verbinden Laufzeitdynamik, Ingenieurregeln,
realisierungsbewusstes Verhalten und Optimierungsnebenbedingungen und
bewahren dabei die Trennung zwischen dauerhaften Alignment-Informationen und
abgeleiteter betrieblicher Interpretation.
\end{summary}
